% Purpose:
% Present baselines as evaluation anchors rather than algorithmic novelties.
%
% Baselines to cover if supported by stable results:
% - no adaptation;
% - feedforward PPO adaptation;
% - pretrained feedforward PPO adaptation;
% - recurrent PPO adaptation.
%
% Tone guidance:
% - Plainly describe what each baseline tests.
% - Do not oversell recurrent PPO as the paper's new method.
% - Do not imply the baseline set is exhaustive.

We use a small baseline family intended to anchor the benchmark rather than to exhaust the ARC adaptation design space. The baselines differ in what information is available before task-specific adaptation and in whether the policy has recurrent state, but they are all evaluated under the same support-to-query protocol.

\paragraph{No adaptation.} The first baseline performs no task-specific optimization. It evaluates the policy as initialized and serves as a lower anchor for how much useful behavior exists before any support-side fitting. In the current package this baseline is most useful on ARC-AGI-1, where it helps separate gains due to actual adaptation from whatever prior bias is already present in the architecture and environment setup.

\paragraph{Feedforward PPO adaptation.} The main baseline performs per-task adaptation with a feedforward PPO policy. This is the simplest direct test of whether standard RL-style inner-loop optimization can use the support pairs to improve downstream behavior. Because it has no recurrent memory, any benefit must come from the task-specific updates themselves rather than from an explicit state that carries support information across episodes.

\paragraph{Pretrained feedforward PPO adaptation.} The pretrained variant starts from a global initialization learned before task-specific adaptation and then applies the same inner-loop PPO update procedure. This baseline asks whether a better starting point improves adaptation efficiency or support-to-query transfer under the same budget ladder. In the current canonical series it does improve support-side similarity on AGI1 relative to the non-pretrained feedforward baseline, which makes it a useful test of whether stronger support fitting alone is enough to move hidden-query performance.

\paragraph{Recurrent PPO adaptation.} The recurrent baseline adds memory to the policy while keeping the basic per-task adaptation framing the same. This baseline is included because ARC tasks naturally raise the question of whether a policy can make better use of support-side structure when it has explicit recurrence. We treat it as a stress test of the benchmark, not as the paper's algorithmic centerpiece.

These four baselines are deliberately modest. They do not cover search-heavy methods, language-model-based adaptation, or broader test-time training schemes that have also become relevant in ARC \citep{cholletARCPrize20242025,coleDontThrowBaby2025}. That limitation is intentional for the first benchmark paper. Our goal is to establish a clean evaluation protocol and show what a stable family of PPO-style baselines can and cannot achieve under it.
